% !TeX encoding = utf-8
% !TeX program = pdflatex

\documentclass{CVM}

\usepackage{amsmath,amssymb}
\usepackage{graphicx}
\usepackage{booktabs}
\usepackage{multirow}
\usepackage{siunitx}
\usepackage{hyperref}

\newcommand{\blu}[1]{\textcolor{blue}{#1}}

\CVMsetup{
type      = {Research Article},
doi       = {CVM.XXXX},
title     = {Cross-Domain Differentiable Monte Carlo Ray Tracing with Embedded
             Physics Proxies for Heliostat Surface Optimization},
author    = {First A. Author$^{1}$, Second B. Author$^{2}$, and Third C. Author$^{3}$\\
\note{}},
runauthor = {F. A. Author, S. B. Author, T. C. Author},
abstract  = {
We present a GPU-native differentiable Monte Carlo ray tracing (DMCRT) system
that embeds a finite-element-derived thin-plate spline (TPS) proxy model and a
20-bin nonlinear gravity deformation library directly into the ray tracing
shaders, enabling the optical loss gradient to back-propagate through both
refraction optics and mechanical deformation simultaneously. The renderer runs
entirely on GPU with zero readback per iteration: a collaborative binary search
computes the S95 spillage threshold on-device, a single Vulkan submit batches
10+ compute dispatches, and only 4 bytes of scalar loss leave the GPU per
iteration. Through a systematic investigation of gravity-induced deformation
mechanisms using this differentiable framework, we design a gradient audit
methodology that decomposes the deformation field into spectral bands, projects
each band onto the bolt influence subspace, and quantifies the structurally
irreducible floor---the fraction of the optical penalty that no bolt adjustment
strategy can remove under a given support layout. Applied to a 20-mirror
heliostat field at Delingha (37.36\si{\degree}N, 97.29\si{\degree}E), the
system reveals that gravity-induced S95 degradation reaches
1.001$\times$--1.539$\times$ with strong spatial non-uniformity predictable
from site geometry alone, and that 84--87\% of the correctable gap is
structurally irreducible under a fixed 35-bolt support layout. Extending the
optimization variable from bolt strokes to the structural layout itself---by
making the support edge margin an optimization variable---yields a 26.7\%
annual S95 reduction with zero additional manufacturing cost. The framework
transforms heliostat surface design from empirical per-mirror adjustment into a
bounded, certifiable, and fully automated differentiable workflow.
},
keywords  = {differentiable rendering, Monte Carlo ray tracing, physics-based
             modeling, GPU computing, heliostat optimization, structural-optical co-design},
copyright = {The Author(s)},
}

% -- address
\address{School of XXX, University XXX, City, Post Code, China. E-mail: corresponding@email.cn.}

\begin{document}

\maketitle

% ===========================================================================
\section{Introduction}
% ===========================================================================

Differentiable rendering has matured rapidly in computer graphics over the
past decade, with systems such as Mitsuba~2/3~\cite{mitsuba2,mitsuba3},
ZeroGrads~\cite{zerograds}, and DiffRP~\cite{diffrp} enabling gradient-based
optimization of scene parameters through physically-based light transport. The
core idea---exposing exact or approximate gradients of pixel values with
respect to scene parameters---has proven transformative for inverse problems in
appearance modeling, geometry reconstruction, and computational
design~\cite{jakob2022dr}.

A parallel development is unfolding in concentrating solar power (CSP)
engineering, where differentiable Monte Carlo ray tracing (DMCRT) has recently
been applied to heliostat aiming optimization~\cite{aiming2025},
canting angle optimization~\cite{canting2025}, and in-situ surface
metrology~\cite{pargmann2024,inverse2025a,inverse2025b}. These works share a
common characteristic: the surface being rendered is \emph{given}---represented
either as a fixed geometric primitive, such as an ideal elliptic paraboloid, or
parameterized by purely optical degrees of freedom (panel tilt angles, focal
lengths). In all cases, the surface shape is treated as a purely geometric
entity. The mechanical degrees of freedom that actually \emph{produce} the
surface shape (support bolt displacements) and the physical loads that
\emph{distort} it (gravity) have never been connected to the optical objective
through a differentiable gradient chain.

% ---------------------------------------------------------------------------
\subsection{Motivation: Why Physics Proxies Matter for Heliostat Optimization}
% ---------------------------------------------------------------------------

The limitations of purely geometric surface representations become evident when
examining the physical reality of a heliostat mirror. A large-span thin glass
mirror ($12.84 \times 9.45$~m, 4~mm thick) supported by discrete bolts sags
under its own weight by several millimeters, and this deformation field changes
continuously as the mirror tracks the sun throughout the day and year. The
resulting slope errors---on the order of several milliradians---are of the same
magnitude as the sun's intrinsic width ($\sim$2.2~mrad for the Buie model at
CSR~0.01), meaning that gravity-induced deformation directly competes with the
optical signal the mirror is designed to concentrate.

Current approaches to modeling surface slope errors in CSP ray tracing fall
into two categories. The first treats slope error as a spatially uniform
Gaussian random variable with a single scalar parameter $\sigma_s$, which
convolves the ideal reflected ray distribution~\cite{hflcal,soltrace}. This
approach is computationally efficient but cannot capture the spatially
structured, tilt-angle-dependent deformation patterns produced by gravity. The
second approach, employed in coupled structural-optical
optimization~\cite{yan2024,yan2025,yang2022,thalange2017}, uses FEA to compute
the actual deformed surface, then passes it to a separate ray tracing
simulation. While physically accurate, the serial, offline nature of this
coupling restricts optimization to parameter sweeps or black-box evolutionary
algorithms over a handful of design variables, with each evaluation requiring a
complete FEA solve plus a complete ray tracing simulation.

Neither approach provides what is needed for gradient-based surface design: a
physics model that is \emph{accurate enough} to capture the dominant
deformation modes, \emph{differentiable} to enable gradient-based optimization,
and \emph{fast enough} to evaluate inside an iterative rendering loop. This is
the gap that an embedded physics proxy aims to fill.

At the same time, embedding physics into a differentiable renderer raises a
deeper question: \emph{how do we know whether the embedded physics is actually
contributing to the optimization, rather than merely existing in the forward
graph?} The question is not academic---the history of physics-informed machine
learning contains cases where a physics term in the loss function or the
forward model was functionally inert, because its gradient contribution was
numerically dominated by other terms or structurally prevented from reaching
the design variables. Answering this question requires dedicated analytical
tools that decompose the physics contribution and quantify what is reachable by
the optimizer and what is not.

% ---------------------------------------------------------------------------
\subsection{Research Questions and Contributions}
% ---------------------------------------------------------------------------

These considerations motivate three research questions that structure this paper:

\begin{itemize}
  \item \textbf{RQ1 (Magnitude):} What is the real annual optical penalty
  caused by gravity-induced deformation, and how does it distribute across a
  heliostat field?

  \item \textbf{RQ2 (Compensability):} For a fixed support structure, what
  fraction of the gravity penalty is \emph{physically removable} by bolt
  adjustment---i.e., where is the floor that separates any reachable design
  from the gravity-free optimum?

  \item \textbf{RQ3 (Structural optimization):} Can we move the floor itself by
  making structural layout parameters---support margin, bolt density, material
  thickness---optimization variables, rather than merely compensating within a
  fixed structure?
\end{itemize}

To answer these questions, we present a complete GPU-native DMCRT system that
embeds a finite-element-derived TPS proxy and a 20-bin NLGEOM gravity library
directly in the ray tracing shaders, enabling end-to-end gradient-based
optimization across both the optical and mechanical domains. Our contributions are:

\begin{enumerate}
  \item \textbf{A zero-readback GPU differentiable rendering pipeline} with
  collaborative on-device S95 threshold search (20 iterations of binary search,
  single workgroup, only \SI{4}{\byte} of scalar loss read back per
  iteration), single-submit batching of 10+ compute dispatches, compile-time
  sunshape specialization, and per-ray pre-culling (\S\ref{sec:system}).

  \item \textbf{A cross-domain gradient chain} that propagates optical loss
  gradients through double-refraction glass optics \emph{and} a
  finite-element-derived thin-plate spline (TPS) mechanical proxy, mapping
  surface flux gradients to 35 independent bolt stroke gradients in a single
  GPU dispatch sequence (\S\ref{sec:proxy}).

  \item \textbf{A gradient audit methodology} that decomposes the gravity
  deformation field into affine, quadratic, and higher-order bands, projects
  each band onto the bolt influence subspace, and quantifies the
  \emph{structurally irreducible floor}---the fraction of the optical penalty
  that no bolt adjustment strategy can remove under the given support layout
  (\S\ref{sec:audit}).

  \item \textbf{A two-level design path} from surface compensation to
  structural optimization: discovering that 80--90\% of gravity slope energy
  concentrates in the cantilevered edge flaps (not between bolts), deriving a
  margin--span conservation law that predicts an interior optimum for the
  support margin, and achieving a \textbf{26.7\% annual S95 reduction} with
  zero additional manufacturing cost (\S\ref{sec:structural}).

  \item \textbf{An application case study} on a 20-mirror heliostat field at
  Delingha, China, providing the first complete gravity penalty map, a
  geometric tilt-angle distribution model that predicts penalty non-uniformity
  from site latitude and tower geometry alone, and experimental evidence that
  grid topology dominates per-bolt position micro-tuning (\S\ref{sec:experiments}).
\end{enumerate}

% ===========================================================================
\section{Related Work}
% ===========================================================================

% ---------------------------------------------------------------------------
\subsection{Differentiable Rendering Systems}
% ---------------------------------------------------------------------------

The modern differentiable rendering stack was established by Mitsuba~2~\cite{mitsuba2},
which introduced retargetable automatic differentiation through a symbolic
computation graph over physically-based light transport. Mitsuba~3~\cite{mitsuba3}
extended this with GPU-accelerated wavefront differentiable path tracing,
enabling optimization over object geometry, materials, and lighting at
practical scales. ZeroGrads~\cite{zerograds} demonstrated that careful
engineering of the AD overhead can make differentiable rendering competitive
with forward-only rendering for production assets.

Our system occupies a distinct point in the design space. Rather than
implementing a general-purpose differentiable renderer with a symbolic
computation graph, we build a \emph{hardware-specialized} pipeline where the
entire gradient computation is written in Slang~\cite{slang} with explicit
\verb|bwd_diff| annotations, compiled to SPIR-V, and dispatched on a single
Vulkan queue with zero CPU--GPU synchronization within an iteration. This
design choice trades generality for throughput: each iteration processes
multiple sun directions (sampled from an annual distribution) in a single
\verb|vkQueueSubmit| call, with only the final scalar loss value crossing the
PCIe boundary. The collaborative on-device S95 binary search
(\S\ref{sec:s95}) eliminates what would otherwise be the dominant readback
bottleneck.

% ---------------------------------------------------------------------------
\subsection{Differentiable MCRT in Concentrating Solar Power}
% ---------------------------------------------------------------------------

Several recent works have applied differentiable ray tracing to CSP problems.
Pargmann et al.~\cite{pargmann2024} used DMCRT for \emph{in-situ} heliostat
metrology, reconstructing surface error fields from calibration images.
Inverse deep learning ray tracing~\cite{inverse2025a,inverse2025b} predicts
heliostat surfaces from focal spot measurements with sim-to-real transfer.
Canting optimization via differentiable ray tracing~\cite{canting2025}
adjusts facet tilt angles to minimize spillage, and aiming optimization via
DMCRT~\cite{aiming2025} optimizes heliostat pointing strategies at the field
level.

These works validate that DMCRT gradients are well-behaved and useful in CSP
optics. However, they all treat the surface as either a given geometric
primitive or parameterized by purely optical variables (tilt, focal length).
The mechanical actuators that physically produce the surface (support bolts)
and the structural loads that physically deform it (gravity) are absent from
the differentiable graph. Our contribution is to close this loop: bolt
displacement $\rightarrow$ TPS surface deformation $\rightarrow$ ray tracing
$\rightarrow$ S95 loss $\rightarrow$ gradient back to bolts, with gravity
entering both the surface height \emph{and} the normal vectors at every step.

% ---------------------------------------------------------------------------
\subsection{Physics Proxies in Differentiable Pipelines}
% ---------------------------------------------------------------------------

The broader trend of embedding physical models into differentiable optimization
has been termed physics-informed machine learning~\cite{karniadakis2021} and has
seen applications from structural optimization~\cite{diffopt2024} to wind farm
digital twins~\cite{ecmwind2023}. The dominant paradigm trains a learned
surrogate (neural network) to approximate the physical simulator, then
differentiates through the surrogate. Our approach is different: we use an
\emph{analytic, hand-derived} proxy (thin-plate splines with a partition-of-unity
property) whose gradients are exact and whose error budget can be bounded by
direct comparison with finite-element solutions (\S\ref{sec:proxy_fidelity}).
This trades the surrogate's flexibility for the proxy's auditability---a choice
that proves essential for the gradient audit in \S\ref{sec:audit}.

% ---------------------------------------------------------------------------
\subsection{Coupled Structural-Optical Heliostat Optimization}
% ---------------------------------------------------------------------------

In the heliostat engineering literature, connecting structural design to optical
performance has been explored through offline, serial coupling: FEA computes
the deformed surface, then a separate ray tracing pass evaluates the optical
consequence~\cite{yan2024,yan2025,yang2022,thalange2017}. These works
established the important baseline that structural parameters (bolt spacing,
margin, material) significantly affect optical quality, and that the
relationship is non-monotonic (an interior optimum exists for bolt spacing).
However, their serial nature restricts optimization to parameter sweeps or
black-box evolutionary algorithms over a handful of design variables, with each
evaluation requiring a complete FEA solve plus a complete ray tracing
simulation.

Our framework differs in three respects: (i) the gradient chain makes the
entire system end-to-end differentiable, enabling gradient-based search over
continuous parameter spaces; (ii) the on-device GPU architecture makes each
evaluation cheap enough for optimization loops (minutes, not hours); and
(iii) the differentiability enables \emph{gradient-reachable subspace
analysis}---answering not just ``which parameter combination is better'' but
``what is the physical upper bound of any adjustment strategy under this
layout'' (\S\ref{sec:audit}). Table~\ref{tab:positioning} positions our work
against the closest prior art.

\begin{table*}[t]
\centering
\caption{Positioning against closest prior work.}
\label{tab:positioning}
\small
\begin{tabular}{lp{2.2cm}p{2.5cm}p{2.2cm}p{1.6cm}p{1.8cm}}
\toprule
Work & Task & Design variables & Structure model & E2E diff. & Bounds analysis \\
\midrule
\cite{aiming2025} & Aiming opt. & 2 DOF/mirror & None & Yes & No \\
\cite{canting2025} & Canting opt. & Panel tilts & None & Yes & No \\
\cite{pargmann2024} & In-situ metrology & Surface normals & None & Yes & No \\
\cite{inverse2025a, inverse2025b} & Surface prediction & Surface coeffs. & None & Neural & No \\
\cite{yan2024,yan2025} & Support $\rightarrow$ optics & Spacing $d$, count $N$ & FEA (offline) & No & No \\
\cite{yang2022,thalange2017} & Structural-optical & Structural dims. & FEA + ray trace & No & No \\
\midrule
\textbf{This work} & \textbf{Surface design} & \textbf{35 bolt strokes + margin} & \textbf{TPS + 20-bin NLGEOM gravity} & \textbf{Yes} & \textbf{Yes} \\
\bottomrule
\end{tabular}
\end{table*}

% ===========================================================================
\section{System Architecture}
\label{sec:system}
% ===========================================================================

Our differentiable renderer is built on Vulkan compute shaders written in
Slang~\cite{slang}, which provides automatic differentiation via
\verb|[Differentiable]| type annotations and \verb|bwd_diff| reverse-mode
code generation. The key architectural decisions are driven by a single
constraint: \emph{every optimization iteration must process multiple sun
directions at interactive rates, with the full gradient chain crossing both
optical and mechanical domains}.

% ---------------------------------------------------------------------------
\subsection{Single-Submit Multi-Dispatch Pipeline}
% ---------------------------------------------------------------------------

The optimization loop iterates over an annual set of sun directions
$\mathcal{D} = \{(\mathbf{s}_d, \text{DNI}_d)\}_{d=1}^{|\mathcal{D}|}$
(typically 36--334 directions, \S\ref{sec:sundir}). For each iteration, a
\emph{single} \verb|vkQueueSubmit| call batches all dispatches for all sun
directions:

\begin{enumerate}
  \item \textbf{Mechanical forward} (\verb|computeBoltSurface|, 1 workgroup):
  TPS influence superposition + 20-bin gravity interpolation $\rightarrow$
  $32\times32$ surface height field $\mathbf{y}$ and normal field $\mathbf{n}$.

  \item \textbf{Ray tracing forward} (\verb|renderForward|, $k$ tiles
  $\times$ $\sim$3950 groups): Monte Carlo sampling from receiver pixels to
  heliostat grid points, double-refraction through \SI{4}{\milli\meter} glass
  ($n=1.523$), Buie sunshape evaluation, flux accumulation in a $157\times50$
  receiver image.

  \item \textbf{GPU S95 search} (\verb|computeS95FindLevel|, 1 workgroup):
  Collaborative binary search for the flux threshold containing 95\% of total
  energy (\S\ref{sec:s95}).

  \item \textbf{Loss + backward optics} (\verb|computeS95LossBUF| +
  \verb|renderBackwardBolt|): Sigmoid loss gradient w.r.t.\ flux $\rightarrow$
  \verb|bwd_diff| through refraction to surface height/normal gradients
  $\partial L/\partial(y, y_u, y_v)$ per grid point.

  \item \textbf{Mechanical backward} (\verb|reduceSurfaceGradients| +
  \verb|projectBoltGradients|): Cross-group reduction of fixed-point
  gradient tiles $\rightarrow$ projection onto bolt gradients
  $\partial L/\partial h_b$ through TPS influence functions.

  \item \textbf{Adam update} (\verb|boltAdamStep|, 1 workgroup): Tanh-bounded
  parameter update in unbounded $\varepsilon$-space, with optional
  regularization terms (anchor, bend, soft stroke wall).
\end{enumerate}

Table~\ref{tab:dispatch} summarizes the per-iteration GPU dispatches. The key
engineering property is that \emph{no intermediate buffer is read back to the
CPU during the optimization loop}: flux images, S95 thresholds, and surface
gradients all remain device-resident, communicated between dispatches via
pipeline barriers. Only the final scalar loss (4 bytes, fixed-point
$\times 10^3$) and per-mirror S95 values cross the PCIe boundary.

\begin{table}[t]
\centering
\caption{GPU dispatches per sun direction per iteration (single submit).}
\label{tab:dispatch}
\small
\begin{tabular}{lccc}
\toprule
Dispatch & Shader & Grid & Purpose \\
\midrule
\verb|computeBoltSurface| & \verb|bolt_forward| & $(1,1,1)$ & Mech.\ forward \\
\verb|clearFlux| & \verb|forward| & $(1,1,1)$ & Zero flux buffer \\
\verb|renderForward| & \verb|forward| (A2 typed) & $(k,3950,1)$ & MCRT forward \\
\verb|finalizeFlux| & \verb|forward| & $(1,1,1)$ & Flux normalization \\
\verb|computeS95FindLevel| & \verb|s95_gpu| & $(1,1,1)$ & S95 binary search \\
\verb|computeS95LossBUF| & \verb|s95_gpu| & $(10,4,1)$ & Loss + dL/dflux \\
\verb|renderBackwardBolt| & \verb|bolt_backward| (A2) & $(3950,k,1)$ & MCRT backward \\
\verb|reduceSurfaceGradients| & \verb|bolt_backward| & $(1,1,1)$ & Cross-group reduce \\
\verb|projectBoltGradients| & \verb|bolt_backward| & $(1,1,1)$ & $\rightarrow \partial L/\partial h_b$ \\
\verb|boltAdamStep| & \verb|bolt_optimizer| & $(1,1,1)$ & Adam update \\
\bottomrule
\end{tabular}
\end{table}

% ---------------------------------------------------------------------------
\subsection{GPU-Resident S95 Threshold Search}
\label{sec:s95}
% ---------------------------------------------------------------------------

The S95 metric (the area of the smallest contiguous set of receiver pixels
containing 95\% of intercepted energy~\cite{s95definition}) serves as the
annual-averaged optical loss. Computing it requires finding a flux threshold
$\tau^*$ such that $\sum_{p: f_p > \tau^*} f_p = 0.95 \cdot E_{\text{total}}$.

A naive implementation would read back the entire flux image ($157\times50 =
7{,}850$ floats), sort on the CPU, and perform the binary search---adding
\SI{\sim30}{\kilo\byte} of PCIe traffic per sun direction per iteration
($\sim$1.1 MB for 36 suns, $\sim$10 MB for 334 suns, per iteration). We
instead implement the search entirely on the GPU as a single-workgroup
collaborative algorithm:

\begin{enumerate}
  \item \textbf{Energy scan:} Each of 256 threads scans its assigned flux
  range to compute the group-wide total energy $E_{\text{total}}$ and maximum
  flux $f_{\max}$, using wave-level reductions.

  \item \textbf{Binary search (20 iterations):} The workgroup maintains a
  search interval $[f_{\text{lo}}, f_{\text{hi}}]$, initially
  $[0, f_{\max}]$. Each iteration sets the probe threshold
  $f_{\text{mid}} = (f_{\text{lo}} + f_{\text{hi}})/2$, then each thread
  accumulates energy above $f_{\text{mid}}$ within its range. A single
  \verb|WorkGroupAllReduce| determines whether the accumulated fraction
  exceeds 0.95, and the interval is halved accordingly.

  \item \textbf{Result:} After 20 iterations, $\tau^*$ is stored to
  \verb|s95State[0]| (a 4-float device buffer, never read back). The total
  energy is stored to \verb|s95State[2]| for use by the energy efficiency
  term (\S\ref{sec:regularization}).
\end{enumerate}

The 20-iteration depth matches a single-precision float mantissa (23 bits),
ensuring convergence to floating-point precision. The relative error compared
to a CPU reference implementation is $\sim 10^{-6}$, limited only by floating-
point reduction order. An environment variable \verb|BEZIER_S95_GPU=0| can
restore the CPU readback path for A/B validation.

% ---------------------------------------------------------------------------
\subsection{Sunshape Specialization and Ray Culling}
\label{sec:optimizations}
% ---------------------------------------------------------------------------

Two further engineering optimizations are worth noting:

\textbf{A2 Compile-time sunshape dispatch} (\verb|forward.slang|): Three
specialized ray tracing entry points (\verb|renderForwardBuie|,
\verb|renderForwardPillbox|, \verb|renderForwardGaussian|) are produced by
Slang's generic specialization with a \verb|let SUNSHAPE_TYPE| compile-time
constant. The C++ dispatcher selects the matching pipeline at runtime based on
the configuration's \verb|sun_type| field. This eliminates a per-ray branch on
the sunshape model, which accounts for $\sim$5\% of inner-loop instructions.

\textbf{A1 Per-ray pre-culling} (\verb|ray_cull|, default on): Before
generating Box-Muller Gaussian samples for the sunshape convolution, each ray
tests $\cos(\text{reflected\_dir}, \text{sun\_dir}) \geq
\cos(\theta_{\text{support}} + \delta_{\text{margin}})$, where
$\theta_{\text{support}} = 0.0436$~rad (Buie cutoff) and
$\delta_{\text{margin}}$ is a configurable margin (default
\SI{8}{\milli\radian}). Rays failing this test are marked invalid via a
bitmask and skipped in both forward and backward passes. At the North
\SI{300}{\meter} mirror, this reduces total iteration time by \textbf{4.8\%}
with exact bitwise agreement to the unculled reference.

% ---------------------------------------------------------------------------
\subsection{Regularization Suite}
\label{sec:regularization}
% ---------------------------------------------------------------------------

The total loss function combines the S95 sigmoid loss with optional
regularization terms, all implemented with analytic gradients in the GPU Adam
shader:

\begin{equation}
\mathcal{L}(\mathbf{h}) = \underbrace{\frac{1}{|\mathcal{D}|}\sum_{d} \text{S95}(\mathbf{h}; \mathbf{s}_d)}_{\text{S95 sigmoid}}
+ \lambda_E \underbrace{\left(\frac{E_{\text{ref}}}{E}-1\right)^2}_{\text{energy efficiency}}
+ \lambda_s \underbrace{(\mathbf{h}-\mathbf{h}^*)^\top \mathbf{G} (\mathbf{h}-\mathbf{h}^*)}_{\text{slope-space anchor}}
+ \lambda_b \underbrace{\mathbf{h}^\top \mathbf{K} \mathbf{h}}_{\text{bending energy}}
+ \lambda_h \underbrace{\sum_b \max(|h_b|-h_{\max},\,0)^2}_{\text{soft stroke wall}}
\label{eq:loss}
\end{equation}

The bolt heights are parameterized as $h_b = h_{\max} \tanh(\varepsilon_b)$
(tanh-bounded, L4), with Adam updates performed in the unbounded
$\varepsilon$-space. The chain-rule factor
$dh_b/d\varepsilon_b = h_{\max}(1-\tanh^2\varepsilon_b)$ is folded
into the gradient, and the learning rate is compensated as
$\text{lr}_\varepsilon = \text{lr} / h_{\max}$ so that the physical step size
near $\varepsilon=0$ (where $h\approx 0$) matches the lr in both bounded and
unbounded modes.

The energy efficiency term $\lambda_E$ prevents the optimizer from ``cheating''
at large slant ranges ($\ge\SI{900}{\meter}$) by spilling most rays off the
receiver to achieve a spuriously small S95. Setting $\lambda_E \ge 0.5$
eliminates this behavior while preserving genuine surface improvements
at short ranges ($\le\SI{600}{\meter}$).

% ===========================================================================
\section{Physics Proxy and Cross-Domain Gradient Chain}
\label{sec:proxy}
% ===========================================================================

The surface model in the plate-local coordinate frame (plate normal = $y$,
plate plane = $x$--$z$) decomposes the normal displacement field into a
gravity term and a bolt superposition term:

\begin{equation}
w(\mathbf{r}; \theta, \mathbf{h}) = w_{\text{grav}}(\mathbf{r}; \theta)
+ \sum_{b=1}^{N_b} h_b \cdot \phi_b(\mathbf{r})
\label{eq:surface}
\end{equation}

where $w_{\text{grav}}$ is the zero-bolt pure-gravity displacement from
NLGEOM finite-element analysis (20 tilt-angle bins, bilinearly interpolated),
and $\phi_b$ is the unit-displacement influence function of bolt $b$,
constructed via thin-plate spline (TPS) interpolation.

% ---------------------------------------------------------------------------
\subsection{TPS Influence Functions}
% ---------------------------------------------------------------------------

For each bolt $b$, we solve the TPS linear system:

\begin{equation}
\begin{bmatrix} \mathbf{K} & \mathbf{P} \\ \mathbf{P}^\top & \mathbf{0} \end{bmatrix}
\begin{bmatrix} \mathbf{c}_b \\ \mathbf{d}_b \end{bmatrix}
= \begin{bmatrix} \mathbf{e}_b \\ \mathbf{0}_3 \end{bmatrix},
\quad K_{ij} = r_{ij}^2 \log(r_{ij}^2),\;
P = [\mathbf{1}, \mathbf{X}, \mathbf{Z}]
\label{eq:tps}
\end{equation}

with Tikhonov regularization $\lambda = 10^{-6}$ on the diagonal
$K_{ii} = \phi(0) + \lambda$. The full-plate response
$\phi_b(x,z) = \sum_j c_{b,j} \cdot r_j^2\log(r_j^2) + d_{b,0} + d_{b,1}x + d_{b,2}z$
satisfies the \textbf{partition of unity} property
$\sum_b \phi_b(\mathbf{r}) \equiv 1$ (peak violation $< 1.3\times 10^{-6}$),
which guarantees physically correct linear superposition of bolt effects.

The system is $38\times38$ for $N_b=35$ bolts ($35 + 3$ polynomial terms),
condition number $\sim 4.2\times 10^6$, solved by dense LU in Python ($< 1$~s).
First-order derivatives are computed analytically:

\begin{equation}
\frac{\partial\phi_b}{\partial u} = \frac{\partial\phi_b}{\partial x}\cdot W,\quad
\frac{\partial\phi_b}{\partial v} = \frac{\partial\phi_b}{\partial z}\cdot L
\label{eq:derivs}
\end{equation}

where $W=12.84$~m, $L=9.45$~m are the plate dimensions, consistent with the
shader's tangent-vector convention $\mathbf{t}_u = (W, y_u, 0)$,
$\mathbf{t}_v = (0, y_v, L)$.

The data are stored as three precomputed $N_b \times 32 \times 32$ float32
buffers, uploaded once to the GPU and accessed via 128 binding slots.

% ---------------------------------------------------------------------------
\subsection{Gravity Library: 20-Bin NLGEOM with Normal Coupling}
% ---------------------------------------------------------------------------

The gravity deformation is obtained from ANSYS MAPDL nonlinear static
simulations (SHELL181 elements, NLGEOM ON, 20 tilt angles
$\theta \in \{10\si{\degree}, 14\si{\degree}, \dots, 80\si{\degree}\}$,
spacing $\le 4\si{\degree}$). The 20-bin density was validated against a
5-bin prototype: the sparse bins produced systematic slope underestimation at
off-bin angles (slope correlation dropping to 0.77--0.82), causing a
``small-spot illusion.'' At 20 bins, the linear interpolation residual is
negligible (R\textsuperscript{2} $\ge 0.96$, slope correlation $\ge 0.994$).

Each bin stores a \textbf{3-plane format}
($w$, $\partial w/\partial u$, $\partial w/\partial v$) on the same
$32\times32$ pixel-centered grid, packed into a single merged buffer of
$20 \times 3 \times 1024$ floats (\SI{240}{\kilo\byte}). The three-plane
design is essential: both the height $w$ and its slopes enter the surface
tangent vectors, so the surface normal becomes

\begin{equation}
\mathbf{n} = -\text{normalize}\big(
\mathbf{t}_u \times \mathbf{t}_v\big),\quad
\mathbf{t}_u = (W,\; y_u^{\text{grav}} + y_u^{\text{bolt}},\; 0),\;
\mathbf{t}_v = (0,\; y_v^{\text{grav}} + y_v^{\text{bolt}},\; L)
\label{eq:normal}
\end{equation}

If only the height were embedded (a 1-plane format), the gravity deformation
would alter the ray intersection point but not the reflected direction---the
gravity effect would be invisible to the optical loss at heliostat-to-receiver
distances. The three-plane embedding ensures that gravity deformation
contributes to the surface normal and therefore to the optical gradient.

% ---------------------------------------------------------------------------
\subsection{Gradient Chain: Optical $\rightarrow$ Mechanical}
% ---------------------------------------------------------------------------

The gradient of the S95 loss with respect to bolt height $h_b$ traces through
three domains:

\begin{equation}
\frac{dL}{dh_b} = \sum_{d\in\mathcal{D}}\sum_{p\in\text{pixels}}
\Bigg[
\frac{\partial L}{\partial f_p} \cdot \Bigg(
\underbrace{\frac{\partial f_p}{\partial y} \cdot \phi_b}_{\text{height channel}}
+ \underbrace{\frac{\partial f_p}{\partial y_u} \cdot \frac{\partial \phi_b}{\partial u}}_{\text{$u$-slope channel}}
+ \underbrace{\frac{\partial f_p}{\partial y_v} \cdot \frac{\partial \phi_b}{\partial v}}_{\text{$v$-slope channel}}
\Bigg)
\Bigg]
\label{eq:gradient}
\end{equation}

The optical partials $\partial f_p / \partial(y, y_u, y_v)$ are computed by
Slang's \verb|bwd_diff| through the double-refraction ray tracing code in
\verb|forward.slang|, accumulated per pixel, and stored in a 12~KB
fixed-point tile ($\times 10^5$ scale, 32-bit unsigned integer, interlocked
atomic addition). A reduction pass (\verb|reduceSurfaceGradients|) converts
the tile to a float3 \verb|surfaceGradient| buffer, and the projection pass
(\verb|projectBoltGradients|) evaluates Eq.~\ref{eq:gradient} using the
precomputed influence buffers.

This three-stage design (per-pixel fixed-point accumulation $\rightarrow$
cross-group float reduction $\rightarrow$ bolt projection) is motivated by a
GPU-specific constraint: \verb|InterlockedAdd| on float is not universally
available (pre-Vulkan~1.2), and a naive global atomic on the bolt gradient
would serialize all thread groups. The fixed-point tile trades one integer
atomic per pixel for one float reduction per sun direction---a \textbf{50$\times$}
reduction in atomic operations for 7,850 pixels and 35 bolts.

% ---------------------------------------------------------------------------
\subsection{Proxy Fidelity}
\label{sec:proxy_fidelity}
% ---------------------------------------------------------------------------

Table~\ref{tab:proxy_fidelity} summarizes the TPS proxy validation against
NLGEOM FEA reference solutions. The systematic deviation (RMS 2.0--3.3~mm,
shape correlation 0.95--0.96) arises because TPS is a purely geometric RBF
interpolation---it does not encode plate bending stiffness $D$, material
properties ($E, \nu$), or free-edge boundary conditions. For the differential
effect $\partial w / \partial h_b$ (which is what the gradient chain uses),
the proxy accuracy is sufficient, as confirmed by the per-pixel flux NRMSE of
1.7--1.8\% and correlation $>0.996$ with FEA ray-traced reference spots.

\begin{table}[t]
\centering
\caption{TPS proxy vs.\ NLGEOM FEA validation (North \SI{300}{\meter}, 35-bolt optimal).}
\label{tab:proxy_fidelity}
\small
\begin{tabular}{lcccc}
\toprule
Angle & RMS (mm) & R\textsuperscript{2} & Shape corr. & Flux corr. \\
\midrule
\SI{29.5}{\degree} & 1.94--1.96 & 0.955 & 0.980 & $>0.997$ \\
\SI{58.5}{\degree} & 2.04--2.05 & 0.952 & 0.980 & $>0.997$ \\
\bottomrule
\end{tabular}
\end{table}

% ===========================================================================
\section{Gradient Audit: Deformation Analysis and Bounds}
\label{sec:audit}
% ===========================================================================

With the physics proxy embedded in the differentiable renderer, we can now use
the framework to answer RQ1 and RQ2: \emph{how large is the gravity penalty,
what are its physical origins, and how much of it can bolt adjustment actually
remove?}

% ---------------------------------------------------------------------------
\subsection{Deformation Decomposition}
% ---------------------------------------------------------------------------

We decompose the per-angle gravity displacement field $w_{\text{grav}}(\mathbf{r}; \theta)$
into three spectral bands via least-squares fitting:

\begin{itemize}
  \item \textbf{Affine} (overall tilt + piston): $w_{\text{affine}} = a_1 x + a_2 z + a_3$.
  \item \textbf{Quadratic} (focus-like): $w_{\text{quad}} = b_1 x^2 + b_2 z^2 + b_3 xz$.
  \item \textbf{Higher-order} (inter-support sag): the residual after removing affine
  and quadratic components.
\end{itemize}

The affine component is \emph{identically zero} at all angles: the 35 support
bolts pin the plate's mean plane, so there is no net tilt or piston to
compensate. This has an important implication: panel-level adjustment methods
(canting, aiming) that act exclusively on affine degrees of freedom cannot,
in principle, touch the dominant gravity damage. The quadratic component is
small (equivalent slope $\le \SI{0.6}{\milli\radian}$). The higher-order
inter-support sag dominates: \textbf{\SI{6.36}{\milli\radian}} equivalent
slope at \SI{10}{\degree} tilt, decreasing to \SI{0.92}{\milli\radian} at
\SI{80}{\degree}.

% ---------------------------------------------------------------------------
\subsection{Compensability: Projection onto the Bolt Subspace}
% ---------------------------------------------------------------------------

To determine how much of this deformation can be compensated by bolt
adjustment, we project the gravity slope field onto the subspace spanned by the
35 TPS influence functions. The projection is performed in \emph{slope space}
because the optical consequence is driven by normal errors, not height errors:

\begin{equation}
\text{removal\_ratio}(\theta) =
\frac{\|\nabla w_{\text{grav}}\|^2 - \min_{\mathbf{h}}\|\nabla w_{\text{grav}} - \sum_b h_b \nabla\phi_b\|^2}
{\|\nabla w_{\text{grav}}\|^2}
\label{eq:removal}
\end{equation}

The result: TPS influence functions can remove only \textbf{26--38\%} of the
slope variance of the higher-order sag component. The remaining 62--74\% is
\textbf{structurally irreducible} under the given 35-bolt (7$\times$5, 8\%
edge margin) support layout.

% ---------------------------------------------------------------------------
\subsection{The Bounds Framework}
% ---------------------------------------------------------------------------

We formalize this into a four-level bounds hierarchy that characterizes the
optimization landscape:

\begin{itemize}
  \item \textbf{$B^*$} (gravity-free optimum): S95 of the best surface
  achievable with gravity disabled---the ultimate lower bound.
  \item \textbf{$B_{\text{ideal}}$} (gravity-free ideal): S95 of the ideal
  elliptic paraboloid surface evaluated without gravity.
  \item \textbf{$B_{\text{comp}}$}: S95 using closed-form slope-space gravity
  compensation as bolt initialization.
  \item \textbf{$B_{\text{naive}}$}: S95 of the ideal elliptic surface
  evaluated with gravity enabled, without any compensation.
\end{itemize}

Table~\ref{tab:bounds} shows these bounds for the four cardinal mirrors at
\SI{300}{\meter} from a \SI{300}{\meter} tower (Delingha, 37.36\si{\degree}N).
The key finding: \textbf{84--87\% of the naive--ideal gap is structurally
irreducible} for the East, South, and West mirrors. Closed-form compensation
recovers only 10--14\%; end-to-end optimization adds at most another 3--5\%.
The floor is a structural property of the support layout, not an artifact of
the optimization setup.

\begin{table}[t]
\centering
\caption{Four-level bounds: S95 (m\textsuperscript{2}) at \SI{300}{\meter},
36 sun directions.}
\label{tab:bounds}
\small
\begin{tabular}{lcccc}
\toprule
Bound & North & East & South & West \\
\midrule
$B^*$ (no-gravity opt.\ lower bound) & 49.77 & 65.00 & 73.07 & 64.68 \\
$B_{\text{ideal}}$ (no-gravity LSQ) & 51.32 & 65.68 & 73.51 & 65.60 \\
$B_{\text{comp}}$ (closed-form comp.)  & 51.75 & 76.71 & 94.90 & 76.83 \\
$B_{\text{naive}}$ (LSQ + gravity)    & 51.31 & 77.90 & 98.33 & 78.07 \\
\midrule
Struct.\ irreducible (\%) & 37\% & \textbf{86.9\%} & \textbf{84.0\%} & \textbf{86.0\%} \\
\bottomrule
\end{tabular}
\end{table}

% ---------------------------------------------------------------------------
\subsection{Ablation Study: The Floor is Invariant}
% ---------------------------------------------------------------------------

A nine-group ablation experiment varied five factors simultaneously:
tanh-bounded vs.\ unbounded parameterization; slope-space anchor strength
$\lambda_s \in \{0, 10^3, 10^4, 10^5\}$; soft stroke wall on/off; bending
energy $\lambda_b$ on/off; and initialization (naive vs.\ compensated). The
result: the East, South, and West mirrors' final S95 values are \textbf{locked
to 76.2, 94.3, and 76.2~m\textsuperscript{2}} respectively, with a range of
$\le 0.13$~m\textsuperscript{2} across all nine groups---within Monte Carlo
noise. The structural floor is invariant to parameterization, regularization
strength, constraint type, and initialization. This invariance constitutes
the primary evidence that the floor is a \emph{structural} property of the
support layout.

% ===========================================================================
\section{From Surface Compensation to Structural Optimization}
\label{sec:structural}
% ===========================================================================

The conclusion that the structural floor is determined by the support layout,
not by bolt stroke optimization, directly motivates making the layout itself an
optimization variable.

% ---------------------------------------------------------------------------
\subsection{The Cantilever Discovery}
% ---------------------------------------------------------------------------

A density scan tested whether adding more bolts within the same 8\% margin
could lower the floor: 48 bolts (8$\times$6), 49 bolts (7$\times$7), 63 bolts
(9$\times$7). The result was unexpected: \textbf{+80\% more bolts (from 35 to
63) reduced the gravity floor by only 4.2\%}. Two competing plate-theory models
(max-span\textsuperscript{4} and Navier simply-supported) both predicted
substantially larger reductions, and were simultaneously falsified.

Spatial partitioning of the slope energy revealed the cause: \textbf{80--91\%}
of the gravity slope energy (across all angles) is concentrated in the
\emph{cantilevered edge flaps}---the region outside the outermost bolt row---not
in the inter-bolt bays that additional bolts would subdivide. The cantilever
is an unsupported free edge, and its deformation scales as $w \propto a^4$
(where $a$ is the overhang distance) in thin-plate theory.

% ---------------------------------------------------------------------------
\subsection{The Margin--Span Conservation Law}
% ---------------------------------------------------------------------------

Reducing the edge margin $m$ (the distance from the outermost bolt row to the
plate edge, as a fraction of the half-width/half-length) simultaneously:

\begin{itemize}
  \item \textbf{Reduces} the cantilever overhang $\propto m$, lowering the
  flap deformation $\propto m^3$;
  \item \textbf{Increases} the inter-bolt span (for a fixed number of bolts),
  raising the bay deformation $\propto (1-m)^3$.
\end{itemize}

The opposing scaling implies an \textbf{interior optimum} for the margin at a
given bolt density. This ``margin--span conservation law'' predicts that
margin optimization can reduce the gravity floor \emph{with zero additional
bolts}---a free structural dividend.

% ---------------------------------------------------------------------------
\subsection{Margin Optimization Results}
% ---------------------------------------------------------------------------

To test this prediction, we scanned the margin across five values
$m \in \{0.08, 0.06, 0.05, 0.04, 0.03\}$ by regenerating the ANSYS gravity
bins at each margin and running the end-to-end optimization pipeline (100
iterations, 110 sun directions, four mirrors). The gravity field for each
margin was obtained from fresh NLGEOM FEA solves with fine substeps
(\verb|NSUBST 50,500,50|) to ensure quasi-static smooth branch following.

The optimum occurs at $m^* = 0.04$--$0.05$, with point estimate $m^* = 0.05$.
At 334 sun directions, the annual-averaged S95 for the four \SI{300}{\meter}
mirrors totals:

\begin{equation}
\text{S95}_{\text{total}}(m{=}0.05) = 281.7\text{ m}^2 \quad\text{vs.}\quad
\text{S95}_{\text{total}}(m{=}0.08) = 384.2\text{ m}^2
\label{eq:final}
\end{equation}

a \textbf{26.7\% reduction} achieved purely by moving the bolt rows closer to
the plate edges---no additional bolts, no material change, no manufacturing
cost increase.

\begin{table}[t]
\centering
\caption{Margin optimization: per-mirror annual S95 (m\textsuperscript{2}),
334 sun directions.}
\label{tab:margin}
\small
\begin{tabular}{lccccc}
\toprule
Margin & North & East & South & West & Total \\
\midrule
0.08 (baseline) & 72.75 & 101.59 & 116.57 & 93.28 & 384.2 \\
0.05 (optimal)  & 54.06 &  73.23 &  81.28 & 73.09 & \textbf{281.7} \\
\midrule
Dividend (\%)    & $-$25.7 & $-$27.9 & $-$30.3 & $-$21.6 & \textbf{$-$26.7} \\
\bottomrule
\end{tabular}
\end{table}

% ---------------------------------------------------------------------------
\subsection{Methodological Note: NLGEOM Substep Requirements}
\label{sec:substep}
% ---------------------------------------------------------------------------

During this investigation, we encountered a systematic artifact in per-angle
independent NLGEOM solves with coarse substeps (ANSYS \verb|NSUBST 1,10,1|).
At tilt angles near \SI{46}{\degree}, the snap-through bistability of the thin
plate causes the independent cold-start solve to randomly select one of two
stable branches. This produces gravity bins with \emph{sign-flipped}
deformation fields at high tilt angles, rendering entire margin-scan results
invalid. The issue was resolved by using fine substeps
(\verb|NSUBST 50,500,50|) to guarantee quasi-static smooth branch following,
which shifted the point estimate from $m^*{=}0.04$ to $m^*{=}0.05$. This
finding constitutes a methodological contribution: \emph{any multi-angle
parameterized model relying on per-angle independent FEA solves must verify
branch continuity with fine substeps or numerical continuation.}

% ===========================================================================
\section{Experiments}
\label{sec:experiments}
% ===========================================================================

% ---------------------------------------------------------------------------
\subsection{Experimental Setup}
% ---------------------------------------------------------------------------

\begin{itemize}
  \item \textbf{Hardware:} NVIDIA RTX 4070 SUPER 12~GB (desktop), RTX 4060
  8~GB (laptop). Vulkan~1.3, Slang compiler.
  \item \textbf{Heliostat:} $12.84 \times 9.45$~m glass mirror, thickness
  \SI{4}{\milli\meter}, $n=1.523$, 35 support bolts (7$\times$5, margin 8\%
  unless otherwise noted). Plate bending stiffness $D=392$~N$\cdot$m, gravity
  load $q=98.1$~N/m\textsuperscript{2}.
  \item \textbf{Receiver:} Cylindrical, radius \SI{10}{\meter}, height
  \SI{20}{\meter}, discretized at $157\times50$ pixels.
  \item \textbf{Sun model:} Buie~\cite{buie2003} CSR$=0.01$, DNI$=1000$~W/m\textsuperscript{2},
  slope error $\sigma_s = 1$~mrad.
  \item \textbf{Sun directions:} 36 directions for fast iteration, 110
  (paper-minimal), 334 (balanced annual), all symmetric about true solar noon
  (azimuth$=180\si{\degree}$, eliminating the equation-of-time asymmetry).
  \item \textbf{Optimizer:} Adam ($\beta_1{=}0.9$, $\beta_2{=}0.999$),
  learning rate $4\times10^{-4}$ constant, 100--200 iterations.
  \item \textbf{Site:} Delingha, Qinghai, China (37.36\si{\degree}N,
  97.29\si{\degree}E), tower height \SI{300}{\meter}.
\end{itemize}

% ---------------------------------------------------------------------------
\subsection{RQ1: The Real Gravity Penalty Map}
\label{sec:exp_rq1}
% ---------------------------------------------------------------------------

Figure~\ref{fig:penalty} shows the annual S95 inflation ratio
$\text{S95}_{\text{naive}} / \text{S95}_{\text{ideal}}$ across 20 mirrors
(five slant ranges 150/300/600/900/1200~m $\times$ four cardinal directions
N/E/S/W, 334 sun directions). The gravity penalty is \textbf{spatially
strongly non-uniform}: North mirrors are nearly immune across all distances
(ratio 1.001--1.02$\times$), while the worst case is the South~\SI{150}{\meter}
mirror at 1.539$\times$.

\begin{figure}[t]
\centering
% \includegraphics[width=\linewidth]{fig_penalty_map}
\caption{Annual gravity penalty map: S95 inflation ratio across 20 mirrors
(334 sun directions). North mirrors are nearly immune; South near-field is
the worst case at 1.539$\times$.}
\label{fig:penalty}
\end{figure}

The non-uniformity has a purely geometric explanation. The annual tilt angle
$\theta$ (angle between plate normal and vertical) distribution for each
mirror position is determined by site latitude, tower geometry, and mirror
coordinates alone---no optical simulation required. At Delingha's latitude:

\begin{itemize}
  \item \textbf{North \SI{300}{\meter}}: $\theta \in [36\si{\degree}, 65\si{\degree}]$,
  median $51\si{\degree}$. 35\% of annual time falls in the 40--50\si{\degree}
  band, near the $\sim$46\si{\degree} NLGEOM zero-crossing where the
  gravity-induced sag is minimal.
  \item \textbf{South \SI{300}{\meter}}: $\theta \in [0\si{\degree}, 57\si{\degree}]$,
  with 24\% of time at $\theta < 20\si{\degree}$, where the gravity bending
  moment $q\cos\theta$ is maximal.
  \item \textbf{East/West \SI{300}{\meter}}: cross the 46\si{\degree} sign-change
  point, carrying an irreducible $\theta$-variation component of
  2.5--2.9~mrad that no fixed bolt setting can compensate.
\end{itemize}

This tilt-angle distribution model predicts the measured penalty map with
Pearson $r = 0.913$ (Spearman $\rho = 0.922$), transforming the empirical
observation into a \textbf{design rule}: which heliostats in a field are worth
structural compensation can be determined from latitude and tower geometry
alone.

% ---------------------------------------------------------------------------
\subsection{RQ2: The Irreducible Floor}
% ---------------------------------------------------------------------------

\S\ref{sec:audit} presented the quantitative bounds. Here we highlight the
experimental evidence for the floor's invariance. The nine-group ablation
(Table~\ref{tab:bounds}) showed that the East, South, and West mirrors' final
S95 is locked to $\sim$76.2, 94.3, and 76.2~m\textsuperscript{2} respectively,
with range $\le 0.13$~m\textsuperscript{2}. Neither the parameterization
(bounded vs.\ unbounded), nor regularization strength (0--$10^5$), nor
initialization (naive vs.\ compensated) could move the East/West mirrors below
76.2 or the South mirror below 94.3.

\textbf{The path forward is structural.} With fixed $N_b = 35$ and fixed
margin $m = 8\%$, bolt stroke adjustment---however sophisticated---can recover
at most $\sim$13--16\% of the gravity penalty. The remaining 84--87\% requires
changing the support layout itself.

% ---------------------------------------------------------------------------
\subsection{RQ3: Structural Optimization}
% ---------------------------------------------------------------------------

\S\ref{sec:structural} detailed the margin optimization. Here we supplement
two additional structural findings.

\textbf{Bolt density at optimal margin.} At the optimal margin $m^*{=}0.05$,
a density scan across layouts 7$\times$5 (35), 9$\times$7 (63), 9$\times$8 (72),
and 11$\times$9 (99 bolts) reveals that 9$\times$7 (63 bolts, +80\%) reduces
S95 by a further 10--14\% beyond the margin dividend. However, 11$\times$9 (99
bolts, +183\%) shows \emph{degradation} on East and West mirrors---the
excessive degrees of freedom cause the optimizer to fall into inferior local
minima. The sweet spot emerges as a ``density $\times$ aspect-ratio'' surface
with 10$\times$7 (70 bolts) producing the overall best results.

\textbf{Position--height joint optimization.} To test whether per-bolt position
freedom beyond uniform grids could yield further gains, we implemented TPS
adjoint position sensitivity and a bold top-K step-ladder protocol. The result:
across four initial grid topologies and two mirror orientations, per-bolt
position moves of only 1--3~cm (from step sizes up to 30~cm offered) produce
S95 improvements of \textbf{0.5--0.9\%} before hitting a wall. \textbf{Grid
topology choice dominates position micro-tuning.} No asymmetric layouts
emerged, confirming that the uniform-grid optimum is near-global under per-bolt
freedom.

% ---------------------------------------------------------------------------
\subsection{Timing and Scalability}
% ---------------------------------------------------------------------------

Table~\ref{tab:timing} reports per-iteration timing. The dominant cost is the
ray tracing forward and backward passes, which scale with pixel count
($157\times50 = 7{,}850$) and sun direction count. The S95 search and
mechanical passes are negligible ($<1$~ms each). A 200-iteration optimization
with 36 sun directions completes in $\sim$7 minutes for a single mirror and
$\sim$20 minutes for four mirrors.

\begin{table}[t]
\centering
\caption{Per-iteration timing breakdown (RTX 4060 Laptop, 36 sun directions,
single mirror).}
\label{tab:timing}
\small
\begin{tabular}{lr}
\toprule
Stage & Time (ms) \\
\midrule
Mechanical forward (surface build) & 0.3 \\
Ray tracing forward (per sun) & 18.2 \\
S95 binary search & 0.1 \\
Loss + ray tracing backward (per sun) & 26.5 \\
Mechanical backward (reduce + project) & 0.5 \\
Adam update & $<0.1$ \\
\midrule
\textbf{Total per sun} & \textbf{45.6} \\
\textbf{Total per iteration (36 suns)} & \textbf{1,640} \\
\bottomrule
\end{tabular}
\end{table}

% ===========================================================================
\section{Discussion}
% ===========================================================================

% ---------------------------------------------------------------------------
\subsection{What Did the Physics Proxy Actually Do?}
% ---------------------------------------------------------------------------

The central question motivating this work---what does embedding a physics proxy
into a differentiable renderer actually achieve?---can now be answered with
quantitative evidence:

\begin{enumerate}
  \item \textbf{Diagnostic role:} The proxy enabled the gradient audit
  methodology that certified the structural floor and proved its invariance to
  all optimizer and regularization choices. The finding that ``84--87\% of the
  gap is structural'' was independently verified by FEA spot-checks.

  \item \textbf{Engineering role:} The proxy carried the deformation field
  through the differentiable loop fast enough to make structural optimization
  practical---replacing per-iteration FEA solves (\SI{3}{\minute} each) with
  GPU-resident $\mu$s evaluations.

  \item \textbf{Explanatory role:} The proxy translated the empirical
  observation ``bolt density shows diminishing returns''~\cite{yan2024,yan2025}
  into a mechanical mechanism---the cantilever-dominant slope energy that the
  margin--span conservation law formalizes.
\end{enumerate}

The proxy's real value lies not in per-mirror compensation (where it can only
recover 13--16\%), but in \emph{driving structural design decisions} (where the
margin dividend of 26.7\% dwarfs the compensation ceiling).

% ---------------------------------------------------------------------------
\subsection{Limitations}
% ---------------------------------------------------------------------------

\textbf{TPS linearity.} The TPS proxy assumes linear superposition of bolt
effects, which the NLGEOM FEA validates with shape correlation 0.95--0.96.
Residual nonlinearity limits the proxy's accuracy for absolute deformation
prediction but is adequate for the differential effect $\partial w/\partial
h_b$ used in gradient computation.

\textbf{Single-site, single-load.} Our experiments consider one site
(Delingha) and one dominant load (gravity). Wind-induced deformation, thermal
gradients, and manufacturing tolerances are not modeled. Incorporating these
loads into the differentiable pipeline would require additional FEA
simulations but no architectural changes.

\textbf{Steel vs.\ glass.} Real heliostats use steel-backed glass modules
(Kesheng~V3: \SI{5}{\milli\meter} steel, rail spacing \SI{0.37}{\meter}).
We verified via FEA that steel and glass produce near-identical deformation
shapes (cosine similarity $>0.996$) with a pure-scaling factor $\alpha \approx
1.01$, meaning the structural conclusions (margin optimum, density surface)
transfer qualitatively. However, the rail-stiffened floor is $\sim$500$\times$
stiffer than the bare plate, so the absolute S95 numbers require recalibration
for production hardware.

% ---------------------------------------------------------------------------
\subsection{Extensions}
% ---------------------------------------------------------------------------

\textbf{Full-field scaling.} The per-mirror S95 results can be converted to
annual field-level energy yield by convolving with the mirror population
distribution and a receiver efficiency model.

\textbf{Aiming co-optimization.} Joint optimization of bolt strokes and
per-mirror aiming points (2 additional DOF per mirror) would close the last
open loop in the optical chain.

\textbf{Branch engineering.} The NLGEOM bistability at $\sim$46\si{\degree}
(\S\ref{sec:substep}) is currently treated as an artifact to be avoided.
However, the flipped branch is \emph{optically better} (lower gravity sag) at
high tilt angles. Exploiting this bistability through controlled snap-through
actuation is a speculative but potentially high-impact extension.

% ===========================================================================
\section{Conclusion}
% ===========================================================================

We have presented a GPU-native differentiable Monte Carlo ray tracing system
that embeds finite-element-derived physics proxies to enable end-to-end
gradient-based optimization of heliostat surfaces across structural and optical
domains. Three main findings structure the work:

\begin{enumerate}
  \item \textbf{The gravity penalty is real, large, and predictable.}
  Gravity-induced S95 inflation reaches 1.539$\times$ in the worst case and
  is strongly non-uniform across the field. The non-uniformity can be
  predicted from site geometry alone, transforming an empirical observation
  into a design rule.

  \item \textbf{The structural floor can be bounded and certified.} Under a
  fixed 35-bolt layout, only 26--38\% of the gravity slope variance is
  reachable by bolt adjustment; 84--87\% of the optical penalty is
  structurally irreducible. This floor is invariant to optimizer,
  regularization, and initialization---a gradient audit methodology with
  applicability to physics-embedded differentiable systems.

  \item \textbf{The floor can be moved by structural optimization.}
  Discovering that 80--90\% of gravity energy concentrates in cantilevered
  edges leads to a margin--span conservation law and a zero-cost 26.7\%
  S95 reduction through margin optimization alone, verified by full ANSYS
  truth re-solve.
\end{enumerate}

The framework transforms gravity compensation from an empirical,
per-mirror adjustment process into a bounded, certifiable, and fully automated
design workflow. The gradient audit methodology---decomposing a physics proxy's
contribution into reachable and irreducible subspaces---is not specific to
heliostats and may serve as a template for other domains where physical
proxies are embedded in differentiable rendering pipelines.

\appendix

\subsection*{Appendix A: ANSYS Automation and Coordinate Conventions}

The ANSYS MAPDL batch pipeline automates per-angle NLGEOM static solves with
SHELL181 shell elements and 35 bolt-node translation constraints. The plate
local coordinate system is defined by the plate normal
$(0, \cos\theta, +\sin\theta)$ in global coordinates, with the displacement
extraction $w = u_y \cos\theta + u_z \sin\theta$.

\subsection*{Appendix B: Sun Direction Sampling}

Sun directions are generated symmetric about \emph{true} solar noon
(azimuth $= 180\si{\degree}$), eliminating the equation-of-time asymmetry that
would bias morning/afternoon sampling in mean-solar-noon schemes. Three
presets are available: \verb|paper| (12 months $\times$ 1 day $\times$ 13
hours, $\sim$110 directions), \verb|balanced| (12 $\times$ 3 $\times$ 13,
$\sim$334 directions), and \verb|dense| (12 $\times$ 14 $\times$ 13,
$\sim$1,556 directions).

\subsection*{Availability of data and materials}

The source code and pre-generated proxy data are available at
\url{https://github.com/xxx/bezier_opt_desktop}. The ANSYS FEA pipeline
requires a commercial ANSYS MAPDL license; pre-generated gravity bins are
included in the repository.

\subsection*{Author contributions}

All authors contributed to the system design, experiments, and manuscript
preparation. All authors have approved the final manuscript.

\subsection*{Acknowledgements}

This work was supported by [TODO: funding information].

\subsection*{Declaration of competing interest}

The authors have no competing interests to declare that are relevant to the
content of this article.

\begin{thebibliography}{99}

\bibitem{mitsuba2}
Nimier-David, M.; Vicini, D.; Zeltner, T.; Jakob, W.
Mitsuba 2: A retargetable forward and inverse renderer.
\textit{ACM Transactions on Graphics} Vol. 38, No. 6, Article No. 203, 2019.

\bibitem{mitsuba3}
Jakob, W.; Speierer, S.; Roussel, N.; Nimier-David, M.; Vicini, D.;
Zeltner, T.; Nicolet, B.; Crespo, M.; Leroy, V.; Zhang, Z.
Mitsuba 3 renderer, 2022. Available at \url{https://mitsuba-renderer.org}.

\bibitem{zerograds}
Yan, Y.; Zhang, Z.; Jakob, W.
ZeroGrads: Learned optimization of production assets via differentiable
rendering. \textit{ACM Transactions on Graphics} Vol. 43, No. 4, 2024.

\bibitem{diffrp}
Zhou, K.; Chen, W.; Yu, J.
Computational caustic design for surface light sources.
\textit{IEEE Transactions on Visualization and Computer Graphics}
Vol. 32, No. 2, 2026.

\bibitem{jakob2022dr}
Jakob, W.
Differentiable rendering. In: \textit{ACM SIGGRAPH 2022 Courses}, 2022.

\bibitem{slang}
He, Y.; Foley, T.; Tatarchuk, N.; Fatahalian, K.
Slang: A language for real-time GPU shading and differentiation.
\textit{ACM Transactions on Graphics} Vol. 41, No. 6, 2023.

\bibitem{aiming2025}
(TBC) A novel heliostat aiming optimization framework via differentiable
Monte Carlo ray tracing for solar power tower.
\textit{Applied Energy}, 2025.

\bibitem{canting2025}
(TBC) The optimization of heliostat paraboloid canting via differentiable
ray tracing. \textit{Solar Energy}, 2025.

\bibitem{pargmann2024}
Pargmann, M.; Ertl, J.; Kesselheim, S.; G\"ottsche, J.; Schorn, C.;
Klink, S.
Automatic heliostat learning for in situ concentrating solar power plant
metrology with differentiable ray tracing.
\textit{Nature Communications} Vol. 15, Article No. 6997, 2024.

\bibitem{inverse2025a}
(TBC) Inverse deep learning raytracing for heliostat surface prediction.
\textit{Solar Energy}, 2025.

\bibitem{inverse2025b}
(TBC) Scalable heliostat surface predictions from focal spots: Sim-to-Real
transfer of inverse deep learning raytracing.
\textit{Solar Energy}, 2025.

\bibitem{karniadakis2021}
Karniadakis, G. E.; Kevrekidis, I. G.; Lu, L.; Perdikaris, P.;
Wang, S.; Yang, L.
Physics-informed machine learning.
\textit{Nature Reviews Physics} Vol. 3, 422--440, 2021.

\bibitem{diffopt2024}
(TBC) Differentiable automatic structural optimization using graph deep
learning. \textit{Advanced Engineering Informatics}, 2024.

\bibitem{ecmwind2023}
(TBC) Digital twin of wind farms via physics-informed deep learning.
\textit{Energy Conversion and Management}, 2023.

\bibitem{yan2024}
Li, B.; Yan, J.; Zhou, W.; Peng, Y.
Influence of service load and structural parameters on optical accuracy
of solar tower heliostat (in Chinese).
\textit{Acta Optica Sinica} Vol. 44, No. 6, Article No. 0623001, 2024.

\bibitem{yan2025}
Yan, J.; Song, T.; Peng, Y.; Zhou, W.
Optical accuracy study of pentagonal tower solar heliostat based on
support bolt jacking molding (in Chinese).
\textit{Acta Optica Sinica} Vol. 45, No. 6, Article No. 0608001, 2025.

\bibitem{yang2022}
Yang, (TBC) et al.
A coupled structural-optical analysis of a novel umbrella heliostat.
\textit{Solar Energy}, 2022.

\bibitem{thalange2017}
Thalange, (TBC) et al.
Design, optimization and optical performance study of tripod heliostat
for solar power tower plant.
\textit{Energy} Vol. 135, 610--624, 2017.

\bibitem{buie2003}
Buie, D.; Monger, A. G.; Dey, C. J.
Sunshape distributions for terrestrial solar simulations.
\textit{Solar Energy} Vol. 74, No. 2, 113--122, 2003.

\bibitem{s95definition}
Collado, F. J.; Guallar, J.
A review of optimized design layouts for solar power tower plants with
Campo code. \textit{Renewable and Sustainable Energy Reviews}
Vol. 20, 142--154, 2013.

\bibitem{fastncm2026}
Lin, R.; (TBC) et al.
Real-time radiative flux density distribution simulation via data-driven
neural convolution model (Fast-NCM). Preprint, 2026.

\bibitem{hflcal}
Ho, C. K.; Khalsa, S. S.
A flux density distribution model based on convolution of Gaussian and
error functions (HFLCAL). \textit{Journal of Solar Energy Engineering}
Vol. 134, No. 1, 2012.

\bibitem{soltrace}
Wendelin, T.
SolTRACE: A new optical modeling tool for concentrating solar optics.
In: \textit{Proceedings of the ASME International Solar Energy Conference},
2003.

\bibitem{tonatiuh}
Blanco, M. J.; et al.
Tonatiuh: An open source Monte Carlo ray tracer for solar concentrating
systems. \textit{Solar Energy}, 2010.

\bibitem{kesheng}
(TBC) Kesheng V3 heliostat specifications. Private communication, 2026.

\end{thebibliography}

\subsection*{Author biography}

\begin{biography}[author1]{First A. Author}
[TODO: author photo and biography]
\end{biography}

\end{document}
